\section{Discussion: each axis bites at a different stage}\label{sec:discussion}

Table~\ref{tab:staged} re-estimates every outcome on one harmonized
sample --- one row per debut owner, scored on that owner's first filing
--- under a single specification, so the coefficients can be read down a
column. Magnitudes differ from the section-specific tables above.

\begin{table}[h]\centering\small
\caption{Every outcome in the funnel, conditioned on its prerequisite, under one specification. Unit is the debut owner, scored on that owner's first filing. Each row is a linear probability model with class$\times$debut-year fixed effects and heteroskedasticity-robust errors; coefficients are percentage points per standard deviation. \emph{Atypicality} is the average of the two KL levels; \emph{lead} is the signed difference. Cohort windows differ because each maintenance gate needs elapsed time and Form~D is electronic only from 2009, so the rows are not a nested decomposition. The first gate is a dated cancellation for non-use; the second is whether a registration that cleared the first was renewed rather than cancelled or expired, restricted to 2002--2013 cohorts, since a year-six and a year-ten death share a terminal status and are separable only this way.}
\label{tab:staged}
\begin{tabular}{llrrrrr}
\toprule
 & & & & \multicolumn{2}{c}{Unsigned} & Signed \\
\cmidrule(lr){5-6}\cmidrule(lr){7-7}
Outcome & Given & $n$ & Base & Atypicality & $|\Delta$KL$|$ & Lead \\
\midrule
Reached registration & filed & 3,161,403 & 56.12\% & $-0.903$ (0.032) & $-0.344$ (0.047) & $+1.095$ (0.029) \\
Passed first gate (yr 6) & registered & 1,774,111 & 47.94\% & $+0.858$ (0.043) & $-0.449$ (0.062) & $-0.955$ (0.038) \\
Passed second gate (yr 10) & passed first & 298,002 & 55.92\% & $+0.705$ (0.102) & $-0.294$ (0.171) & $-0.861$ (0.104) \\
Raised a Reg D round & filed & 1,534,824 & 2.53\% & $-0.023$ (0.015) & $+0.034$ (0.030) & $-0.272$ (0.018) \\
Raised a Reg D round & registered & 916,442 & 2.61\% & $-0.056$ (0.020) & $-0.015$ (0.040) & $-0.235$ (0.023) \\
Listed after the round & funded & 38,886 & 3.49\% & $+0.445$ (0.123) & $+0.618$ (0.208) & $-0.641$ (0.128) \\
Reached listed universe & registered & 1,774,111 & 0.23\% & $+0.030$ (0.005) & $+0.008$ (0.008) & $-0.002$ (0.005) \\
\bottomrule
\end{tabular}
\\[0.3em]\footnotesize Standard errors in parentheses, in place of significance stars; see the note to Table~\ref{tab:gate_lpm}.
\end{table}

Surprising language costs at registration and pays afterwards: atypical
debut filings complete registration less often ($-1.35$pp per standard
deviation), then survive the five-year proof of continued use more often
($+1.35$pp) and reach public markets more often ($+0.04$pp on a 0.56\%
base --- the one specification in which a positive atypicality
coefficient on listing survives, and it is a continuous term through a
non-monotone profile). Having been early runs the other way: leading
applications complete registration \emph{more} often ($+0.87$pp) and then
pass the five-year proof \emph{less} often ($-0.49$pp on passing). Both
survival rows point the same way at the year-ten renewal, four years
later and on a sixth of the sample: atypicality $+1.30$pp and lead
$-0.63$pp among marks that had already cleared the five-year proof.
Whatever the five-year proof selects on, the renewal selects on it again
rather than reversing it.

In the financing rows atypicality does nothing at one stage and a great
deal at the next. Whether a debut owner ever raises a Regulation~D round
is unrelated to how atypical its filing was ($-0.017$pp per standard
deviation against a 2.61\% base, $t = -0.8$); among owners who did raise
one, the same measure predicts \emph{reaching} the public market
($+0.49$pp on a 3.49\% base, $t = 3.8$). Atypical language is invisible
at the financing stage and informative afterwards, which is what an
investor screen that does not read the text would look like --- the text
being free and public at the moment of filing. The comparison is among
survivors, so it cannot separate that reading from selection on
unobservables. Three limits bound these rows: Regulation~D notices are
electronic only from 2009; owners are matched by normalized name, which
succeeds for a minority favouring formally constituted entities; and the
final row conditions on funding realized after the filing being scored,
so it establishes that information distinguishing eventual listers was
present in the text years earlier, not that the language caused the
outcome.

Lead does not reverse anywhere after registration: negative at both use
proofs, at financing with or without registration imposed, and at listing
among the funded; the only row where it is not negative is SEC reporting
among registered owners, where it is indistinguishable from zero. Leading
filings fare worse at every stage where the market, rather than the
examiner, does the selecting.

Why would having been early cost? The two axes are not independent in the
obvious economic sense. What makes a filing leading is precisely that its
industry moved toward its language, and an industry that has moved toward
your language now describes its goods the way you describe yours. Being
early therefore consumes the quantity that pays: the mark that was
unusual in 2011 is ordinary by 2016 not because its owner changed
anything but because the category arrived. On this reading atypicality is
a position competitors erode, and lead measures how fast. The design
cannot demonstrate that mechanism; what can be said is that a negative
coefficient on lead at equal atypicality is what the mechanism would
produce if it were real.

\subsection{Implications for trademark-based measurement}\label{sec:implications}

Three practices follow for research that reads trademark text. First,
score distributions, not words. Word-level divergence is dominated by
description length --- on this corpus the correlation with log distinct
terms reaches $-0.68$ --- and firm-level performance correlations built
on it dissolve or reverse under distribution scoring on identical
samples. Any term-level novelty measure should be accompanied by the
theme-level rescoring that costs one model fit. Second, treat the
registered description as a conformed document. Applicants who refile
rewrite toward their class's typical language from both tails, whoever
drafts the second attempt, so samples of registered marks understate the
atypicality applicants first attempted, and estimates on registered text
are estimates on prosecution-filtered language. Studies using description
similarity or breadth inherit this compression. Third, the two components
of surprise are not interchangeable and should not be pooled. Lead is
robust to every partition, weighting, and window tried here;
atypicality's apparent effects can be artifacts of the partition (its
protective effect at the five-year proof belongs to theme importation,
visible only to global models) or of a control (its listing effect
depends on how description length is held). A study that reports
``novelty'' from trademark text without separating the temporal from the
cross-sectional component may attribute to one what belongs to the other.

For users of the classification itself, one caution: internet-bearing
filings in goods classes survive at rates closer to the technology
classes than to their own, so Nice classes used as industry controls do
not partition the internet economy. And for the maintenance data, one
asset: the five-year proof of continued use is a dated, sworn,
product-level survival record covering every registered mark --- an
outcome the indicator literature has largely left unused, observable for
any cohort more than seven years old at no cost beyond parsing the
prosecution history.
